\documentclass[11pt,a4paper,oneside]{report}
\usepackage[spanish,activeacute,es-noshorthands]{babel}
\usepackage[utf8]{inputenc}
\usepackage[T1]{fontenc}
\usepackage{lmodern}
\usepackage{amsmath, amssymb}
\usepackage{graphicx}
\usepackage[dvipsnames]{xcolor}
\usepackage{geometry}
\geometry{margin=2.4cm, headheight=14pt}
\usepackage{microtype}
\usepackage{booktabs}
\usepackage{longtable}
\usepackage{enumitem}
\usepackage{fancyhdr}
\usepackage{tcolorbox}
\tcbuselibrary{skins, breakable, listings}
\usepackage{caption}
\captionsetup{font=small, labelfont=bf}
\usepackage{titlesec}
\usepackage{listings}
\usepackage{hyperref}
\hypersetup{
  colorlinks=true,
  linkcolor=NavyBlue,
  citecolor=ForestGreen,
  urlcolor=BrickRed,
  pdftitle={Raspberry Pi 5 --- Funcionamiento profundo},
  pdfauthor={CubeSat-EdgeAI-Payload}
}

\definecolor{intisat}{RGB}{26, 54, 93}
\definecolor{intisataccent}{RGB}{197, 143, 0}
\definecolor{ok}{RGB}{47, 133, 90}
\definecolor{warn}{RGB}{197, 60, 60}
\definecolor{codebg}{RGB}{248, 249, 250}

\lstdefinestyle{shell}{
  basicstyle=\ttfamily\footnotesize,
  backgroundcolor=\color{codebg},
  frame=single, rulecolor=\color{gray!30},
  framesep=4pt, breaklines=true,
  showstringspaces=false,
  literate={a}{{a}}1 {e}{{e}}1 {i}{{i}}1 {o}{{o}}1 {u}{{u}}1 {n}{{n}}1
}
\lstset{style=shell}

\titleformat{\chapter}[hang]
  {\normalfont\Huge\bfseries\color{intisat}}
  {\thechapter.}{0.6em}{}
\titlespacing*{\chapter}{0pt}{-20pt}{20pt}

\pagestyle{fancy}
\fancyhf{}
\fancyhead[L]{\nouppercase{\leftmark}}
\fancyhead[R]{\thepage}
\renewcommand{\headrulewidth}{0.4pt}

\newtcolorbox{nota}[1][]{
  colback=blue!4, colframe=intisat,
  fonttitle=\bfseries, title=Para tu payload,
  breakable, sharp corners=south, #1
}
\newtcolorbox{cuidado}[1][]{
  colback=red!4, colframe=warn,
  fonttitle=\bfseries, title=Cuidado,
  breakable, sharp corners=south, #1
}
\newtcolorbox{tip}[1][]{
  colback=green!5, colframe=ok,
  fonttitle=\bfseries, title=Tip practico,
  breakable, sharp corners=south, #1
}

\begin{document}

% ─── Portada ────────────────────────────────────────────────────────
\begin{titlepage}
\centering
\vspace*{1cm}
{\color{intisat}\rule{\textwidth}{2pt}}
\vspace{0.6cm}

{\Huge\bfseries\color{intisat} Raspberry Pi 5}\\[0.3em]
{\LARGE\bfseries Funcionamiento profundo}

\vspace{0.7cm}
{\Large\itshape Hardware, software y subsistemas\\
para el payload CubeSat-EdgeAI}

\vspace{1cm}
{\color{intisataccent}\rule{0.6\textwidth}{1pt}}

\vspace{2cm}

\begin{tabular}{rl}
\textbf{Plataforma:} & Raspberry Pi 5 (4 GB) --- desarrollo \\
\textbf{Modelo de vuelo:} & Compute Module 5 (proximamente) \\
\textbf{Proyecto:} & CubeSat-EdgeAI-Payload \\
\textbf{Repositorio:} & \texttt{github.com/JaminYC/CubeSat-EdgeAI-Payload} \\
\textbf{Fecha:} & \today \\
\textbf{Documento:} & v1.0
\end{tabular}

\vfill

\begin{minipage}{0.85\textwidth}
\textbf{Resumen.}\quad
Documento de referencia integral del Raspberry Pi 5 enfocado en su uso
como payload de procesamiento autonomo a bordo de un CubeSat. Cubre el
SoC BCM2712, la arquitectura ARMv8, el sistema de memoria LPDDR4X, los
buses GPIO (I2C, SPI, UART), el subsistema de camara CSI-2, la PMIC
RP1, el proceso de boot, el sistema operativo Raspberry Pi OS, la
gestion termica, y la migracion al Compute Module 5 para el modelo de
vuelo. Cada capitulo incluye notas especificas para integracion al
INTISAT, comandos practicos de Linux y advertencias de hardware.
\end{minipage}

\vspace{1.2cm}
{\color{intisat}\rule{\textwidth}{2pt}}
\end{titlepage}

\tableofcontents
\clearpage

% =====================================================================
\chapter{Vision general}

El Raspberry Pi 5 es una computadora single-board de credit-card size que
combina un SoC moderno (BCM2712), 4--16 GB de RAM, conectividad
multiproposito y un IO controller dedicado (RP1). Para tu payload actua
como \emph{computadora principal del experimento}: corre Linux, controla
los sensores opticos por buses estandar, ejecuta el pipeline de IA y
dialoga con el OBC del satelite por I2C.

\section{Arquitectura de bloques}

\begin{center}
\begin{tabular}{|p{4cm}|p{10cm}|}
\hline
\textbf{SoC BCM2712} & 4 nucleos ARM Cortex-A76 @ 2.4 GHz, GPU VideoCore VII, ISP, decoder H.265 \\
\hline
\textbf{RP1 IO controller} & chip dedicado (Raspberry Pi diseñado) que expone GPIO, I2C, SPI, UART, USB, Ethernet \\
\hline
\textbf{LPDDR4X RAM} & 4/8/16 GB, bus 64 bit, 4267 MT/s \\
\hline
\textbf{PCIe Gen 2 x1} & FPC connector externo, hasta 5 Gbit/s --- para NVMe SSD \\
\hline
\textbf{2x USB 3.0 + 2x USB 2.0} & via RP1 |
\\
\hline
\textbf{2x camera CSI-2 + 2x display DSI} & pueden funcionar en simultaneo \\
\hline
\textbf{40-pin GPIO header} & compatible con Pi anterior pero con drivers nuevos en RP1 \\
\hline
\end{tabular}
\end{center}

\section{Diferencia clave vs Pi 4}

El cambio mas importante: \textbf{el GPIO ya no esta directamente en el SoC}.
Pasa por el RP1, que es un chip independiente con su propio firmware. Esto
trae implicaciones:

\begin{itemize}[leftmargin=2em]
  \item Latencia GPIO mas alta ($\sim 1\;\mu$s vs ns en Pi 4).
  \item \texttt{RPi.GPIO} y \texttt{wiringpi} no funcionan; hay que usar \texttt{lgpio} o \texttt{gpiozero}.
  \item \texttt{pigpio} todavia funciona (necesario para tu I2C esclavo).
  \item Frecuencias maximas: I2C hasta 1 MHz, SPI hasta 50 MHz.
\end{itemize}

\begin{nota}
El \texttt{cubesat/i2c\_slave.py} usa \texttt{pigpio.bsc\_i2c()} que funciona
en Pi 5 con el demonio pigpiod corriendo. No usa el bus I2C clasico (que
solo soporta master), sino el BSC (Broadcom Serial Controller) que si
soporta modo slave.
\end{nota}

% =====================================================================
\chapter{El SoC BCM2712}

\section{CPU --- ARM Cortex-A76}

Cuatro nucleos identicos, cada uno con:

\begin{itemize}[leftmargin=2em]
  \item Pipeline out-of-order de 12 etapas, 4 instrucciones decodificadas por ciclo.
  \item Cache L1 de instrucciones $64$ KB + L1 de datos $64$ KB.
  \item Cache L2 privada por nucleo de $512$ KB.
  \item Cache L3 compartida de $2$ MB.
  \item Instrucciones SIMD: NEON (Advanced SIMD) y SVE (Scalable Vector Extension).
  \item Frecuencia base $2.4$ GHz, dynamic frequency scaling de $1.5$ a $2.4$ GHz.
\end{itemize}

\textbf{Performance practico} (single-core, comparativa con Pi 4):

\begin{center}
\begin{tabular}{lcc}
\toprule
\textbf{Benchmark} & \textbf{Pi 4} & \textbf{Pi 5} \\
\midrule
Sysbench CPU prime  & 1.5x ref  & 4.2x ref  \\
Geekbench 6 single  & 220       & 770       \\
Geekbench 6 multi   & 600       & 1700      \\
\bottomrule
\end{tabular}
\end{center}

Aproximadamente $3\times$ mas rapido que Pi 4 en cargas single-thread.

\section{GPU --- VideoCore VII}

\begin{itemize}[leftmargin=2em]
  \item 800 MHz, soporta OpenGL ES 3.1 y Vulkan 1.2.
  \item Decoder H.265 hardware (4kp60), decoder H.264 (1080p60), encoder H.264 (1080p30).
  \item NO soporta CUDA (no es NVIDIA).
\end{itemize}

\textbf{Para tu payload}: la GPU NO se usa para IA (no hay framework optimizado
ARM Mali/VideoCore para PyTorch). El pipeline IA corre en CPU.

\section{ISP (Image Signal Processor)}

Procesa la salida cruda de la camara CSI-2 antes de pasarla a memoria:

\begin{itemize}[leftmargin=2em]
  \item Bayer demosaicing.
  \item Auto white balance (AWB).
  \item Auto exposure (AE).
  \item Lens shading correction.
  \item Noise reduction (NR).
\end{itemize}

\begin{cuidado}
El ISP \emph{aplica correcciones automaticas} que pueden interferir con
microscopia cuantitativa. En \texttt{cubesat/capture.py} se desactivan
con \texttt{Picamera2.set\_controls(\{"AwbEnable": False, "AeEnable": False\})}.
\end{cuidado}

\section{Memoria LPDDR4X}

\begin{itemize}[leftmargin=2em]
  \item 4267 MT/s (efectivo $\sim 17$ GB/s en burst).
  \item Bus 64 bit (4 canales x16 internos).
  \item Latencia $\sim 80$ ns (timing CL40).
  \item Versiones: $4, 8, 16$ GB. La de tu payload es $4$ GB.
\end{itemize}

% =====================================================================
\chapter{El controlador RP1 (clave para entender el Pi 5)}

\section{Que es}

El RP1 es un chip diseñado por Raspberry Pi Ltd (no Broadcom) que actua como
\emph{southbridge}: maneja todos los IO perifericos. Esta interconectado al
SoC por PCIe Gen 2 x4.

\section{Que controla}

\begin{itemize}[leftmargin=2em]
  \item 40-pin GPIO header (todos los pines).
  \item 2 buses I2C (I2C1 y dedicated).
  \item 5 buses SPI.
  \item 6 UARTs.
  \item Ethernet Gigabit.
  \item 2x USB 3.0 + 2x USB 2.0.
  \item Audio analog out, PWM, etc.
\end{itemize}

\section{Por que importa para vos}

a) \textbf{Latencia}: el viaje SoC $\to$ PCIe $\to$ RP1 $\to$ pin agrega $\sim 1\;\mu$s.
   No es problema para microscopia, pero seria limitante para señales en
   tiempo real $> 100$ kHz.

b) \textbf{Drivers nuevos}: si una libreria asume que el GPIO esta en el SoC
   (como \texttt{RPi.GPIO}), no funciona. Usar \texttt{lgpio}, \texttt{gpiozero}
   o \texttt{pigpio}.

c) \textbf{I2C slave}: solo el BSC original del SoC soporta modo slave; el
   RP1 expone solo modo master. El \texttt{cubesat/i2c\_slave.py} usa el BSC
   directamente con \texttt{pigpio.bsc\_i2c()}.

% =====================================================================
\chapter{Sistema de potencia}

\section{Rails de tension}

El Pi 5 entra con $5$ V (USB-C o pin 2/4 del header) y la PMIC interna genera:

\begin{center}
\begin{tabular}{lll}
\toprule
\textbf{Rail} & \textbf{Tension} & \textbf{Para que} \\
\midrule
$V_{\text{IN}}$        & 5.0 V & entrada externa \\
$V_{\text{CORE}}$      & 0.85--0.9 V & nucleos CPU (variable, DVFS) \\
$V_{\text{GPU}}$       & 0.9 V       & GPU \\
$V_{\text{DDR}}$       & 1.1 V       & RAM LPDDR4X \\
$V_{\text{IO\_18}}$    & 1.8 V       & buses internos del SoC \\
$V_{\text{IO\_33}}$    & 3.3 V       & GPIO + perifericos externos \\
\bottomrule
\end{tabular}
\end{center}

\section{Consumo (medido)}

\begin{center}
\begin{tabular}{lcc}
\toprule
\textbf{Estado} & \textbf{Corriente @ 5V (A)} & \textbf{Potencia (W)} \\
\midrule
Boot peak               & 1.6 A & 8.0 \\
Idle (HDMI desconectado) & 0.6 A & 3.0 \\
CPU 100\% (1 core)      & 0.9 A & 4.5 \\
CPU 100\% (4 cores)     & 1.4 A & 7.0 \\
CPU + GPU + WiFi + USB  & 1.6 A & 8.0 \\
\bottomrule
\end{tabular}
\end{center}

Mediciones tipicas; la POE+ HAT oficial provee hasta $5$ A.

\section{Brown-out detection}

Si $V_{\text{IN}}$ cae bajo $4.6$ V el Pi entra en \emph{undervoltage}, se
ve un icono de rayo en pantalla y la CPU baja a $1.5$ GHz. En log de kernel:

\begin{lstlisting}[language=bash]
[123.456] Under-voltage detected! (0x00050005)
\end{lstlisting}

Verificar con \texttt{vcgencmd get\_throttled}: el bit 0 indica undervoltage.

\begin{cuidado}
En la integracion al INTISAT, el rail $5$V\_PAYLOAD debe sostener $1.6$ A
de pico sin caer abajo de $4.7$ V. El fusible PTC de $5$ A elegido tiene
margen, pero verificar regulacion del LDO bajo carga.
\end{cuidado}

% =====================================================================
\chapter{El header de 40 pines (GPIO)}

\section{Pinout completo}

\begin{center}
\footnotesize
\begin{tabular}{|c|l|l|c||c|l|l|c|}
\hline
\textbf{Pin} & \textbf{Funcion} & \textbf{BCM} & \textbf{Tension} & \textbf{Pin} & \textbf{Funcion} & \textbf{BCM} & \textbf{Tension} \\
\hline
1  & 3.3V        & ---    & 3.3V & 2  & 5V          & ---    & 5V \\
3  & SDA1 / I2C  & GPIO2  & 3.3V & 4  & 5V          & ---    & 5V \\
5  & SCL1 / I2C  & GPIO3  & 3.3V & 6  & GND         & ---    & 0V \\
7  & GPIO4       & GPIO4  & 3.3V & 8  & TXD0 / UART & GPIO14 & 3.3V \\
9  & GND         & ---    & 0V   & 10 & RXD0 / UART & GPIO15 & 3.3V \\
11 & GPIO17      & GPIO17 & 3.3V & 12 & GPIO18      & GPIO18 & 3.3V \\
13 & GPIO27      & GPIO27 & 3.3V & 14 & GND         & ---    & 0V \\
15 & GPIO22      & GPIO22 & 3.3V & 16 & GPIO23      & GPIO23 & 3.3V \\
17 & 3.3V        & ---    & 3.3V & 18 & GPIO24      & GPIO24 & 3.3V \\
19 & MOSI / SPI  & GPIO10 & 3.3V & 20 & GND         & ---    & 0V \\
21 & MISO / SPI  & GPIO9  & 3.3V & 22 & GPIO25      & GPIO25 & 3.3V \\
23 & SCLK / SPI  & GPIO11 & 3.3V & 24 & CE0 / SPI   & GPIO8  & 3.3V \\
25 & GND         & ---    & 0V   & 26 & CE1 / SPI   & GPIO7  & 3.3V \\
27 & ID\_SDA     & GPIO0  & 3.3V & 28 & ID\_SCL     & GPIO1  & 3.3V \\
29 & GPIO5       & GPIO5  & 3.3V & 30 & GND         & ---    & 0V \\
31 & GPIO6       & GPIO6  & 3.3V & 32 & GPIO12      & GPIO12 & 3.3V \\
33 & GPIO13      & GPIO13 & 3.3V & 34 & GND         & ---    & 0V \\
35 & GPIO19      & GPIO19 & 3.3V & 36 & GPIO16      & GPIO16 & 3.3V \\
37 & GPIO26      & GPIO26 & 3.3V & 38 & GPIO20      & GPIO20 & 3.3V \\
39 & GND         & ---    & 0V   & 40 & GPIO21      & GPIO21 & 3.3V \\
\hline
\end{tabular}
\end{center}

\section{Pines usados por tu payload}

\begin{center}
\begin{tabular}{lll}
\toprule
\textbf{Pin} & \textbf{Funcion} & \textbf{Conectado a} \\
\midrule
2 / 4   & 5V IN              & EPS via fusible + LDO \\
6, 14, 20, 30, 39 & GND      & GND comun \\
3 (GPIO2) & SDA I2C\_2        & TXS0108E $\to$ OBC INTISAT \\
5 (GPIO3) & SCL I2C\_2        & TXS0108E $\to$ OBC INTISAT \\
12 (GPIO18) & OLED DC         & SSD1351 \\
18 (GPIO24) & OLED RST        & SSD1351 \\
19 (GPIO10) & MOSI SPI0       & SSD1351 \\
23 (GPIO11) & SCLK SPI0       & SSD1351 \\
24 (GPIO8)  & CE0 SPI0        & SSD1351 \\
\bottomrule
\end{tabular}
\end{center}

\section{Funciones alternativas (alt modes)}

Cada GPIO puede configurarse en hasta 6 funciones distintas. Por ejemplo
GPIO12 puede ser GPIO digital, PWM, SPI3 MOSI, UART4 TX, etc. Se selecciona
con \texttt{raspi-gpio set <pin> a<n>} o desde \texttt{config.txt}.

\begin{lstlisting}[language=bash]
# Ver configuracion actual
raspi-gpio get 12
# Cambiar a PWM
raspi-gpio set 12 a0
\end{lstlisting}

% =====================================================================
\chapter{Buses de comunicacion}

\section{I2C}

Bus de 2 hilos (SDA, SCL) con multimaster. Velocidades estandar $100$ kHz,
$400$ kHz, $1$ MHz, $3.4$ MHz (high-speed mode no soportado).

\subsection{Direcciones}

I2C usa direcciones de 7 bit. El payload ocupa \textbf{0x42}. Direcciones
reservadas:

\begin{itemize}[leftmargin=2em]
  \item \texttt{0x00}: general call.
  \item \texttt{0x08-0x0F}: high-speed master.
  \item \texttt{0x78-0x7F}: 10-bit addressing.
\end{itemize}

\subsection{Modo slave (BSC)}

El SoC tiene un periferico llamado BSC (Broadcom Serial Controller) que
puede actuar como slave. Solo accesible via \texttt{pigpio.bsc\_i2c()},
no por la libreria estandar.

\begin{lstlisting}[language=python]
import pigpio
pi = pigpio.pi()
pi.bsc_i2c(0x42)   # responde como slave en 0x42

while True:
    s, b, d = pi.bsc_i2c(0x42)   # lee bytes recibidos
    if b > 0:
        process(d[:b])
\end{lstlisting}

\subsection{Test de comandos basicos}

\begin{lstlisting}[language=bash]
# Detectar dispositivos en el bus
sudo i2cdetect -y 1

# Leer 1 byte de address 0x42, registro 0
i2cget -y 1 0x42 0x00

# Escribir 0x05 al registro 0x10 del slave 0x42
i2cset -y 1 0x42 0x10 0x05
\end{lstlisting}

\section{SPI}

Bus sincrono full-duplex. El Pi 5 tiene 5 buses SPI (SPI0-SPI5) via RP1,
mas SPI6 directo del SoC.

\begin{itemize}[leftmargin=2em]
  \item Velocidades: $0.5$ MHz a $50$ MHz.
  \item Modos: 0, 1, 2, 3 (combinaciones de CPOL y CPHA).
  \item En tu payload: SPI0 con OLED SSD1351 a 1 MHz.
\end{itemize}

Habilitar via \texttt{raspi-config} o en \texttt{/boot/firmware/config.txt}:
\begin{lstlisting}[language=bash]
dtparam=spi=on
\end{lstlisting}

\section{UART}

Hasta 6 UARTs disponibles (UART0-UART5). El UART0 esta conectado por defecto
a la consola serial (115200, 8N1). Util para debug:

\begin{lstlisting}[language=bash]
# Conectarse desde otra Pi por UART (pin 8 TX, pin 10 RX)
sudo screen /dev/ttyAMA0 115200
\end{lstlisting}

\section{CSI-2 (camara)}

MIPI Camera Serial Interface 2. El Pi 5 tiene 2 conectores FPC de 22 pines
(diferente del Pi 4 que era de 15 pines, requiere cable adaptador).

\begin{itemize}[leftmargin=2em]
  \item 4 lanes de datos diferenciales.
  \item Hasta 1.5 Gbps por lane.
  \item Soporta hasta $4032 \times 3040$ @ 30 fps (12 MP).
\end{itemize}

\textbf{Para OV5647}: 2 lanes, $2592 \times 1944$ @ 15 fps en 5 MP. Usa
\texttt{Picamera2} library:

\begin{lstlisting}[language=python]
from picamera2 import Picamera2
cam = Picamera2()
config = cam.create_still_configuration(main={"size": (2592, 1944)})
cam.configure(config)
cam.start()
img = cam.capture_array()
\end{lstlisting}

\section{PCIe}

PCIe Gen 2 x1 expuesto en el conector FPC lateral. Velocidad $5$ Gbps. Los
HATs oficiales permiten conectar SSD NVMe.

\begin{nota}
Para el modelo de vuelo CM5, considerar agregar SSD NVMe en lugar de SD
para almacenar resultados. La SD se corrompe facilmente con vibracion y
radiacion. NVMe + eMMC (ya integrada en el CM5) es mucho mas robusta.
\end{nota}

% =====================================================================
\chapter{Almacenamiento}

\section{Tarjeta micro-SD}

Slot estandar SDIO 3.0, soporta hasta UHS-I (104 MB/s). Es el medio principal
en Pi 5 retail.

\textbf{Formato recomendado}:
\begin{itemize}[leftmargin=2em]
  \item Tabla de particiones MBR (no GPT).
  \item Particion 1: FAT32 \texttt{boot} ($\geq 256$ MB).
  \item Particion 2: ext4 \texttt{rootfs} (resto).
\end{itemize}

\begin{cuidado}
Las SD baratas se degradan rapido con escrituras. Para el payload de vuelo
NO USAR SD. Migrar a CM5 con eMMC integrada.
\end{cuidado}

\section{NVMe (via PCIe HAT)}

Velocidades reales $\sim 800$ MB/s read, $700$ MB/s write con un Samsung 980.
Booteo desde NVMe disponible en EEPROM $\geq 2024-09$.

\section{eMMC (solo CM5)}

El Compute Module 5 trae eMMC integrada de $0$ (Lite), $4$, $16$, $32$, $64$ GB.
Mas robusta que SD frente a vibracion y temperatura. Velocidad $\sim 200$ MB/s.

% =====================================================================
\chapter{Proceso de boot}

\section{Etapas del boot}

\begin{enumerate}[leftmargin=2em]
  \item \textbf{Bootloader EEPROM} (interno al SoC): inicializa el SoC,
        habilita LPDDR4X, busca medio de boot segun \texttt{BOOT\_ORDER}.
  \item \textbf{start.elf} (en partition de boot): carga el firmware del GPU.
  \item \textbf{config.txt}: archivo de configuracion (resoluciones, overclocks,
        device tree overlays).
  \item \textbf{cmdline.txt}: argumentos del kernel.
  \item \textbf{Kernel Linux} (\texttt{kernel\_2712.img}): carga el sistema.
  \item \textbf{initramfs}: filesystem temporal con drivers.
  \item \textbf{systemd}: lanza servicios de userspace.
  \item \textbf{Login}: sistema listo.
\end{enumerate}

\section{config.txt importante para tu payload}

\begin{lstlisting}[language=bash]
# Habilitar buses
dtparam=i2c_arm=on
dtparam=spi=on
dtparam=audio=off                    # OLED SSD1351 reusa pines audio

# Camara CSI-2
camera_auto_detect=1
dtoverlay=ov5647

# Watchdog hardware (cuelgue $\to$ reboot automatico)
dtparam=watchdog=on

# I2C slave en GPIO2/3
dtoverlay=i2c-gpio,bus=10,i2c_gpio_sda=2,i2c_gpio_scl=3
\end{lstlisting}

\section{Boot order}

Configurable en EEPROM. Por defecto:
\begin{lstlisting}[language=bash]
BOOT_ORDER=0xf41   # SD, NVMe, USB, repeat
\end{lstlisting}

Editar con \texttt{sudo rpi-eeprom-config -e}.

\section{Recovery / fastboot}

Si el sistema no bootea, mantener pulsado el boton del Pi 5 al darle power
para entrar a un mini OS de recovery. Solo en Pi 5, no Pi 4.

% =====================================================================
\chapter{Sistema operativo}

\section{Raspberry Pi OS}

Distribucion oficial basada en Debian. Version actual al momento de redactar:
\textbf{Bookworm} (Debian 12). Imagenes:

\begin{itemize}[leftmargin=2em]
  \item \textbf{Lite}: sin escritorio, ideal para servidor / payload (lo que usamos).
  \item \textbf{Desktop}: con LXDE, ~5 GB.
  \item \textbf{Full}: con LibreOffice, Chromium, etc., ~10 GB.
\end{itemize}

Bajar de \url{https://www.raspberrypi.com/software/operating-systems/}.

\section{Kernel}

Linux $\sim 6.6$ con patches RPi. Verificar:

\begin{lstlisting}[language=bash]
uname -a
# Linux raspberrypi 6.6.20+rpt-rpi-2712 #1 SMP PREEMPT Debian
\end{lstlisting}

\textbf{PREEMPT} indica preemptive scheduling (mejor latencia).

\section{Systemd}

Sistema de init y manager de servicios. Comandos basicos:

\begin{lstlisting}[language=bash]
# Listar servicios activos
systemctl list-units --type=service --state=running

# Habilitar / deshabilitar
sudo systemctl enable pigpiod
sudo systemctl start pigpiod

# Ver estado
systemctl status cubesat-pipeline

# Logs
journalctl -u cubesat-pipeline -f --since "1 hour ago"
\end{lstlisting}

\section{Servicios del payload}

\begin{itemize}[leftmargin=2em]
  \item \texttt{pigpiod.service}: demonio pigpio (necesario para I2C slave).
  \item \texttt{cubesat-pipeline.service}: el daemon principal con WatchdogSec.
  \item \texttt{cubesat-i2c-slave.service}: dispatcher I2C en 0x42.
  \item \texttt{cubesat-telemetry.service}: snapshots de telemetria cada 30 s.
\end{itemize}

% =====================================================================
\chapter{Programacion de GPIO en Pi 5}

\section{Librerias disponibles}

\begin{center}
\begin{tabular}{lll}
\toprule
\textbf{Libreria} & \textbf{Pi 4} & \textbf{Pi 5} \\
\midrule
\texttt{RPi.GPIO}     & si  & \textbf{NO} (deprecada) \\
\texttt{wiringpi}     & si  & \textbf{NO} (no portada) \\
\texttt{lgpio}        & si  & \textbf{si} (recomendada) \\
\texttt{gpiozero}     & si  & \textbf{si} (alto nivel) \\
\texttt{pigpio}       & si  & \textbf{si} (con demonio) \\
\texttt{libgpiod}     & si  & \textbf{si} (la oficial Linux) \\
\bottomrule
\end{tabular}
\end{center}

\section{Ejemplo: blink de LED en GPIO17}

\begin{lstlisting}[language=python]
import gpiozero
import time
led = gpiozero.LED(17)
while True:
    led.on()
    time.sleep(0.5)
    led.off()
    time.sleep(0.5)
\end{lstlisting}

\section{Ejemplo: leer sensor con I2C}

\begin{lstlisting}[language=python]
import smbus2
bus = smbus2.SMBus(1)         # I2C-1
addr = 0x68                   # MPU-6050 IMU
data = bus.read_byte_data(addr, 0x75)   # WHO_AM_I
print(hex(data))              # debe imprimir 0x68
\end{lstlisting}

% =====================================================================
\chapter{Camara y video}

\section{libcamera + Picamera2}

\texttt{libcamera} es el stack moderno de camaras de Linux, reemplaza al viejo
\texttt{raspistill}. \texttt{Picamera2} es el wrapper Python.

\section{Configuracion para microscopia cuantitativa}

\begin{lstlisting}[language=python]
from picamera2 import Picamera2
import numpy as np

cam = Picamera2()
config = cam.create_still_configuration(
    main={"size": (2592, 1944), "format": "RGB888"},
    raw={"size": (2592, 1944)},
)
cam.configure(config)

# Desactivar correcciones automaticas
cam.set_controls({
    "AwbEnable": False,
    "AeEnable": False,
    "ExposureTime": 50000,    # 50 ms
    "AnalogueGain": 1.0,
    "ColourGains": (1.0, 1.0),
})
cam.start()

# Capturar
img = cam.capture_array("main")     # numpy array (H, W, 3)
raw = cam.capture_array("raw")      # Bayer 10-bit
\end{lstlisting}

\section{Sensor OV5647 specs}

\begin{itemize}[leftmargin=2em]
  \item Resolucion activa: $2592 \times 1944$ ($5$ MP).
  \item Pixel pitch: $1.4\;\mu$m.
  \item Profundidad: $10$ bits crudos (RAW10).
  \item Frame rate: $30$ fps @ $1080$p, $15$ fps @ $5$ MP.
  \item Lente original: removida para uso lensless.
\end{itemize}

% =====================================================================
\chapter{Gestion termica}

\section{Sensores}

El SoC BCM2712 tiene un sensor de temperatura interno accesible por:

\begin{lstlisting}[language=bash]
vcgencmd measure_temp
# temp=45.2'C

cat /sys/class/thermal/thermal_zone0/temp
# 45200 (mili-grados Celsius)
\end{lstlisting}

\section{Throttling}

Cuando $T > 80$°C el SoC reduce frecuencia automaticamente:

\begin{itemize}[leftmargin=2em]
  \item $T > 80$°C: throttle suave ($-200$ MHz).
  \item $T > 85$°C: throttle agresivo (1.5 GHz).
  \item $T > 90$°C: posible apagado por proteccion.
\end{itemize}

Verificar con:
\begin{lstlisting}[language=bash]
vcgencmd get_throttled
# throttled=0x0    (sin throttling)
# throttled=0x40000  (en throttling actual)
\end{lstlisting}

\section{Disipacion en CubeSat}

\begin{cuidado}
El CubeSat NO tiene aire para conveccion. La disipacion del Pi 5 (hasta $7$
W en CPU 100\%) genera $\Delta T$ alto. Soluciones:
\begin{itemize}[leftmargin=1em]
  \item Disipador termico de aluminio en contacto directo con el SoC.
  \item Pad termico o pasta hacia la estructura del satelite (que actua como
        radiador).
  \item Limitar el duty cycle del modo PROCESSING.
  \item Considerar undervoltaje y limite de frecuencia: \\
        \texttt{arm\_freq=1500} en \texttt{config.txt} reduce a $1.5$ GHz.
\end{itemize}
\end{cuidado}

% =====================================================================
\chapter{Networking (en lab, no en orbita)}

\section{Ethernet Gigabit}

PHY integrada en RP1, conector RJ-45. Soporta PoE+ (con HAT oficial).

\section{WiFi y Bluetooth}

Cypress CYW43455 dual-band (2.4 + 5 GHz) WiFi 802.11ac + Bluetooth 5.0 BLE.
En CubeSat no se usa (pero util para debug en lab):

\begin{lstlisting}[language=bash]
# Configurar WiFi sin GUI
sudo nmtui
\end{lstlisting}

\section{En el modelo de vuelo}

CM5 \textbf{NO incluye} WiFi/BT por defecto (hay version con). Para
CubeSat se desactivan via \texttt{config.txt}:

\begin{lstlisting}[language=bash]
dtoverlay=disable-wifi
dtoverlay=disable-bt
\end{lstlisting}

% =====================================================================
\chapter{Migracion al Compute Module 5}

\section{Diferencias clave}

\begin{center}
\begin{tabular}{lll}
\toprule
\textbf{Aspecto} & \textbf{Pi 5} & \textbf{CM5} \\
\midrule
Tamaño              & $85 \times 56$ mm     & $55 \times 40$ mm \\
Conector            & headers (40-pin etc.) & mezzanine 200-pin \\
Carrier             & propia placa          & se diseña aparte \\
USB                 & 4 puertos             & no expuestos \\
Ethernet            & RJ-45                 & solo PHY (necesita conector externo) \\
HDMI                & si                    & exposed via mezzanine \\
eMMC                & no                    & \textbf{si (4--64 GB)} \\
Consumo             & 3--8 W                & 2.5--7 W ($\sim 10\%$ menos) \\
SD                  & si                    & solo en version Lite \\
\bottomrule
\end{tabular}
\end{center}

\section{Carrier board para CubeSat}

El carrier debe incluir:

\begin{enumerate}[leftmargin=2em]
  \item Conector mezzanine de 200 pines (Hirose DF40 stack height $1.5$ mm).
  \item Footprint mecanico PC-104.
  \item Level shifter TXS0108E para I2C externo.
  \item Fusible PTC + LDO para el rail $5$ V de la EPS.
  \item INA219 para monitor de potencia.
  \item Conector CSI-2 (15 o 22 pines) para OV5647.
  \item Conector SPI + GPIO para OLED SSD1351.
  \item LED de estado de boot.
  \item Pads de prueba para los buses.
\end{enumerate}

\section{Software}

El CM5 corre el mismo Raspberry Pi OS que el Pi 5. La transicion es
transparente para el codigo: solo cambia el carrier y la imagen de boot
(habilitar eMMC en lugar de SD).

\begin{nota}
La migracion al CM5 NO requiere cambios en \texttt{cubesat/} ni en
\texttt{pipeline/}. El unico cambio sera en \texttt{cubesat/install.sh}
para flashear la eMMC en lugar de la SD.
\end{nota}

% =====================================================================
\chapter{Robustez para vuelo}

\section{Watchdog hardware}

El BCM2712 tiene un watchdog interno que reinicia el SoC si no recibe kick.
Habilitarlo:

\begin{lstlisting}[language=bash]
# config.txt
dtparam=watchdog=on

# /etc/systemd/system/cubesat-pipeline.service
WatchdogSec=600
\end{lstlisting}

Si el daemon se cuelga $\geq 600$ s sin llamar \texttt{sd\_notify(WATCHDOG=1)},
systemd lo reinicia. Si systemd se cuelga, el watchdog hardware reinicia el SoC.

\section{Memoria SD vs eMMC en orbita}

\begin{center}
\begin{tabular}{lll}
\toprule
\textbf{Aspecto} & \textbf{SD} & \textbf{eMMC (CM5)} \\
\midrule
Vibracion          & sale del slot (riesgo) & soldada \\
Radiacion          & sin proteccion        & ECC interna \\
Vida util          & 1000-3000 ciclos     & 10000-20000 ciclos \\
Velocidad          & 100 MB/s              & 200 MB/s \\
\bottomrule
\end{tabular}
\end{center}

\section{Bit rot y radiacion}

En LEO se reciben $\sim 5-10$ rad/dia. El BCM2712 no es rad-hard pero:

\begin{itemize}[leftmargin=2em]
  \item TID toleramos $\sim 10$ krad antes de degradacion notable $=$ unos
        3 años de mision.
  \item SEU (Single Event Upset) en RAM: bit-flip aislado, mitigable con ECC
        software o reboots periodicos.
  \item SEFI (Single Event Functional Interrupt): cuelgue del SoC $\to$
        watchdog hardware reinicia.
\end{itemize}

\section{Uptime y reboots planificados}

Reboot semanal preventivo:

\begin{lstlisting}[language=bash]
# /etc/cron.d/weekly-reboot
0 3 * * 0 root /sbin/shutdown -r now
\end{lstlisting}

% =====================================================================
\chapter{Comandos utiles de diagnostico}

\section{Hardware}

\begin{lstlisting}[language=bash]
# Modelo y revision
cat /proc/device-tree/model
# Raspberry Pi 5 Model B Rev 1.0

# Frecuencia actual de cada core
for i in 0 1 2 3; do
    cat /sys/devices/system/cpu/cpu$i/cpufreq/scaling_cur_freq
done

# Memoria
free -m

# Temperatura
vcgencmd measure_temp

# Tensiones
for v in core sdram_c sdram_i sdram_p; do
    echo -n "$v: "
    vcgencmd measure_volts $v
done

# Tirar el dump completo
vcgencmd version
vcgencmd get_throttled
\end{lstlisting}

\section{Buses}

\begin{lstlisting}[language=bash]
# I2C scan
sudo i2cdetect -y 1

# SPI test
sudo spi-pipe -d /dev/spidev0.0 -s 1000000

# UART loopback
echo "test" | sudo socat - /dev/ttyAMA0,b115200,raw
\end{lstlisting}

\section{Camara}

\begin{lstlisting}[language=bash]
# Listar camaras detectadas
libcamera-hello --list-cameras

# Captura rapida
libcamera-still -o test.jpg --immediate

# Sin auto-exposure
libcamera-still -o test.jpg --immediate --shutter 50000 --gain 1.0
\end{lstlisting}

\section{Servicios}

\begin{lstlisting}[language=bash]
# Estado de todos los cubesat-*
systemctl status cubesat-* --no-pager

# Log en vivo
journalctl -f -u cubesat-pipeline

# Reiniciar todo
sudo systemctl restart cubesat-pipeline cubesat-i2c-slave cubesat-telemetry

# Ver dependencias
systemctl list-dependencies cubesat-pipeline
\end{lstlisting}

% =====================================================================
\chapter{Diagnosticos comunes y soluciones}

\section{El sistema no botea}

\begin{itemize}[leftmargin=2em]
  \item LED rojo encendido continuo: error de power. Verificar fuente $\geq 5$ A.
  \item LED verde parpadea N veces: codigo de error documentado en
        \url{https://github.com/raspberrypi/firmware/issues}.
  \item Sin LEDs: probable fallo de boot ROM. Conectar UART (pin 8/10) para ver mensajes.
\end{itemize}

\section{I2C no detecta dispositivos}

\begin{enumerate}[leftmargin=2em]
  \item \texttt{lsmod | grep i2c} -- ¿estan los modulos cargados?
  \item \texttt{ls /dev/i2c-*} -- ¿hay nodos?
  \item \texttt{sudo i2cdetect -y 1} -- escanea el bus.
  \item Verificar pull-ups (4.7 k$\Omega$ a 3.3 V; suelen estar en la placa de sensor).
  \item Confirmar tension $V_{IO}$ correcta del slave (3.3 V para Pi GPIO).
\end{enumerate}

\section{Camara no detectada}

\begin{lstlisting}[language=bash]
# Re-habilitar
sudo raspi-config
# Interface Options $\to$ Camera $\to$ Enable

# Verificar overlay
grep -E "camera|ov5647" /boot/firmware/config.txt

# Detectar
libcamera-hello --list-cameras
\end{lstlisting}

\section{Temperatura alta}

\begin{itemize}[leftmargin=2em]
  \item Verificar que el disipador este montado y tenga pasta termica.
  \item En \texttt{config.txt}: \texttt{arm\_freq=1500} para limitar a 1.5 GHz.
  \item Reducir \texttt{over\_voltage=-2} (undervolt seguro).
  \item Verificar duty cycle del pipeline IA (no correr permanentemente).
\end{itemize}

% =====================================================================
\chapter{Checklist final pre-vuelo}

\begin{itemize}[leftmargin=2em]
  \item[$\square$] Migrado a Compute Module 5 con eMMC.
  \item[$\square$] Carrier board diseñado y validado en banco.
  \item[$\square$] Watchdog hardware habilitado (\texttt{dtparam=watchdog=on}).
  \item[$\square$] Watchdog software habilitado (\texttt{WatchdogSec=600}).
  \item[$\square$] Auto-reboot semanal configurado en cron.
  \item[$\square$] WiFi/BT deshabilitados.
  \item[$\square$] HDMI deshabilitado.
  \item[$\square$] Audio deshabilitado.
  \item[$\square$] Frecuencia limitada a $1.5$ GHz.
  \item[$\square$] Disipador con pasta termica + pad a estructura.
  \item[$\square$] OTA testado en banco (git fetch + checkout + rollback).
  \item[$\square$] Burn-in 168 h sin reboot.
  \item[$\square$] Test termico --10 a +50 °C.
  \item[$\square$] Test vibracion 3 g RMS.
  \item[$\square$] Documentacion de mision completa en repo.
\end{itemize}

% =====================================================================
\chapter*{Referencias oficiales}
\addcontentsline{toc}{chapter}{Referencias oficiales}

\begin{itemize}[leftmargin=2em]
  \item Raspberry Pi 5 product page: \url{https://www.raspberrypi.com/products/raspberry-pi-5/}
  \item Pi 5 documentation: \url{https://www.raspberrypi.com/documentation/computers/raspberry-pi-5.html}
  \item BCM2712 datasheet: \url{https://datasheets.raspberrypi.com/bcm2712/bcm2712-peripherals.pdf}
  \item RP1 datasheet: \url{https://datasheets.raspberrypi.com/rp1/rp1-peripherals.pdf}
  \item CM5 datasheet: \url{https://datasheets.raspberrypi.com/cm5/cm5-datasheet.pdf}
  \item CM5 IO Board: \url{https://datasheets.raspberrypi.com/cm5/cm5-ioboard-datasheet.pdf}
  \item Picamera2 manual: \url{https://datasheets.raspberrypi.com/camera/picamera2-manual.pdf}
  \item libcamera apps: \url{https://www.raspberrypi.com/documentation/computers/camera_software.html}
  \item OV5647 datasheet: \url{https://www.ovt.com/sensors/OV5647}
  \item SSD1351 datasheet: \url{https://www.solomon-systech.com/product/ssd1351/}
  \item Forum oficial: \url{https://forums.raspberrypi.com}
  \item GitHub Pi firmware: \url{https://github.com/raspberrypi/firmware}
\end{itemize}

\end{document}
